% ----------------------------------------------------------------------
% Parte II, capítulo 4
% ----------------------------------------------------------------------
\chapter{Percurso metodológico}
\label{cap:metodo}

\instrucao{Descreva o desenho da pesquisa com detalhe suficiente para outro pesquisador repeti-la: universo, amostra, instrumentos, período de coleta, tratamento dos dados e limitações do método.}

\section{Universo e amostra}
O universo compreende os repositórios institucionais brasileiros registrados em diretórios internacionais. A amostra intencional reuniu repositórios de universidades públicas com interface OAI-PMH ativa no período da coleta.

\section{Coleta}
\instrucao{Quando a coleta for automatizada, informe a ferramenta, os parâmetros e as datas. O roteiro detalhado pode ir para um apêndice.}
A coleta teve duas frentes: a colheita automática dos registros pelo protocolo OAI-PMH, descrita no Apêndice~\ref{ap:coleta}, e um questionário às equipes gestoras. A especificação do protocolo usada como referência está resumida no Anexo~\ref{an:oai}.

\section{Indicadores}
Cada dimensão do modelo foi traduzida em indicadores binários ou proporcionais. Por exemplo, a dimensão semântica considera a proporção de registros com autor vinculado a um identificador persistente, como o ORCID.
